% GENERATED FILE. Edit canonical lessons, then run npm run generate:books.
% canonical-source-hash: fnv1a64:ef842f12ed66059b
% canonical-lessons: JA-C214-arashi, JA-C214-yureru, JA-C214-hi-fire, JA-C214-moeru, JA-C214-kieru

\chapter{Storms, Shaking and Fire}
\label{ch:ja-storms-shaking-and-fire}
\hlchaptermodality{214}

\begin{chapteropening}
\textbf{By the end of this chapter:} \emph{I can talk about storms, shaking and fire.}

Complete the last lesson of chapter 214: I can talk about storms, shaking and fire.
\end{chapteropening}

\section[\texorpdfstring{\ja{あらし} (arashi)}{あらし (arashi)}]{\ja{あらし} (arashi) — a storm}

\label{lesson:JA-C214-arashi}



\begin{quote}
Before the new one: say the Japanese for to climb, then the Japanese for to get warm.
\end{quote}

\subsection*{You'll want to know: \ja{あらし}}
\textbf{\ja{あらし}} — \emph{arashi} — \textquotedblleft{}a storm\textquotedblright{}.

\begin{quote}
\textbf{\ja{あしたは} \ja{あらしでしょう。}} — \emph{ashita wa arashi deshō} — \textquotedblleft{}There will probably be a storm tomorrow.\textquotedblright{}
\end{quote}

\begin{cousinweb}
\textbf{\ja{あらし}} (\emph{arashi}) is a native Japanese word (wago), inherited rather than borrowed from Chinese; its kanji was attached later as a kun'yomi, and its older origin is not settled (JMdict, KANJIDIC2).
\end{cousinweb}

\begin{grammarlens}[title={What you've built}]
One more word for everyday life at A2.
\end{grammarlens}

\subsection*{Your turn}
\begin{itemize}
\raggedright
  \item \emph{Say it:} \emph{arashi}
  \item \emph{Say it:} \emph{arashi}, once more
  \item \emph{Say it:} \emph{ashita wa arashi deshō} — There will probably be a storm tomorrow.
  \item \emph{Recall:} say the Japanese for to climb, then the Japanese for to get warm, then say \emph{arashi} again
\end{itemize}

\begin{tcolorbox}[breakable,colback=teal!4,colframe=teal!35!black,title={Before you move on}]
What does \ja{あらし} mean? (\textquotedblleft{}A storm\textquotedblright{}.) Say it once more.
\end{tcolorbox}
\section[\texorpdfstring{\ja{ゆれる} (yureru)}{ゆれる (yureru)}]{\ja{ゆれる} (yureru) — to shake; to sway}

\label{lesson:JA-C214-yureru}



\begin{quote}
Before the new one: say the Japanese for to get warm, then the Japanese for a storm.
\end{quote}

\subsection*{You'll want to know: \ja{ゆれる}}
\textbf{\ja{ゆれる}} — \emph{yureru} — \textquotedblleft{}to shake; to sway\textquotedblright{}.

\begin{quote}
\textbf{\ja{かぜで} \ja{きが} \ja{ゆれます。}} — \emph{kaze de ki ga yuremasu} — \textquotedblleft{}The trees sway in the wind.\textquotedblright{}
\end{quote}

Polite \textbf{\ja{ゆれます}}, \emph{yuremasu}; an ichidan verb.

\begin{cousinweb}
\textbf{\ja{ゆれる}} (\emph{yureru}) is a native Japanese word (wago), inherited rather than borrowed from Chinese; its kanji was attached later as a kun'yomi, and its older origin is not settled (JMdict, KANJIDIC2).
\end{cousinweb}

\begin{grammarlens}[title={What you've built}]
One more word for everyday life at A2.
\end{grammarlens}

\subsection*{Your turn}
\begin{itemize}
\raggedright
  \item \emph{Say it:} \emph{yureru}
  \item \emph{Say it:} \emph{yureru}, once more
  \item \emph{Say it:} \emph{kaze de ki ga yuremasu} — The trees sway in the wind.
  \item \emph{Recall:} say the Japanese for to get warm, then the Japanese for a storm, then say \emph{yureru} again
\end{itemize}

\begin{tcolorbox}[breakable,colback=teal!4,colframe=teal!35!black,title={Before you move on}]
What does \ja{ゆれる} mean? (\textquotedblleft{}To shake; to sway\textquotedblright{}.) Say it once more.
\end{tcolorbox}
\section[\texorpdfstring{\ja{ひ} (hi)}{ひ (hi)}]{\ja{ひ} (hi) — a fire; a flame}

\label{lesson:JA-C214-hi-fire}



\begin{quote}
Before the new one: say the Japanese for a storm, then the Japanese for to shake.
\end{quote}

\subsection*{You'll want to know: \ja{ひ}}
\textbf{\ja{ひ}} — \emph{hi} — \textquotedblleft{}a fire; a flame\textquotedblright{}.

\begin{quote}
\textbf{\ja{ひの} \ja{ちかくは} \ja{あついです。}} — \emph{hi no chikaku wa atsui desu} — \textquotedblleft{}It is hot near the fire.\textquotedblright{}
\end{quote}

\begin{cousinweb}
\textbf{\ja{ひ}} (\emph{hi}) is a native Japanese word (wago), inherited rather than borrowed from Chinese; its kanji was attached later as a kun'yomi, and its older origin is not settled (JMdict, KANJIDIC2).
\end{cousinweb}

\begin{grammarlens}[title={What you've built}]
One more word for everyday life at A2.
\end{grammarlens}

\subsection*{Your turn}
\begin{itemize}
\raggedright
  \item \emph{Say it:} \emph{hi}
  \item \emph{Say it:} \emph{hi}, once more
  \item \emph{Say it:} \emph{hi no chikaku wa atsui desu} — It is hot near the fire.
  \item \emph{Recall:} say the Japanese for a storm, then the Japanese for to shake, then say \emph{hi} again
\end{itemize}

\begin{tcolorbox}[breakable,colback=teal!4,colframe=teal!35!black,title={Before you move on}]
What does \ja{ひ} mean? (\textquotedblleft{}A fire; a flame\textquotedblright{}.) Say it once more.
\end{tcolorbox}
\section[\texorpdfstring{\ja{もえる} (moeru)}{もえる (moeru)}]{\ja{もえる} (moeru) — to burn}

\label{lesson:JA-C214-moeru}



\begin{quote}
Before the new one: say the Japanese for to shake, then the Japanese for a fire.
\end{quote}

\subsection*{You'll want to know: \ja{もえる}}
\textbf{\ja{もえる}} — \emph{moeru} — \textquotedblleft{}to burn\textquotedblright{}.

\begin{quote}
\textbf{\ja{かみは} \ja{よく} \ja{もえます。}} — \emph{kami wa yoku moemasu} — \textquotedblleft{}Paper burns well.\textquotedblright{}
\end{quote}

Polite \textbf{\ja{もえます}}, \emph{moemasu}; an ichidan verb.

\begin{cousinweb}
\textbf{\ja{もえる}} (\emph{moeru}) is a native Japanese word (wago), inherited rather than borrowed from Chinese; its kanji was attached later as a kun'yomi, and its older origin is not settled (JMdict, KANJIDIC2).
\end{cousinweb}

\begin{grammarlens}[title={What you've built}]
One more word for everyday life at A2.
\end{grammarlens}

\subsection*{Your turn}
\begin{itemize}
\raggedright
  \item \emph{Say it:} \emph{moeru}
  \item \emph{Say it:} \emph{moeru}, once more
  \item \emph{Say it:} \emph{kami wa yoku moemasu} — Paper burns well.
  \item \emph{Recall:} say the Japanese for to shake, then the Japanese for a fire, then say \emph{moeru} again
\end{itemize}

\begin{tcolorbox}[breakable,colback=teal!4,colframe=teal!35!black,title={Before you move on}]
What does \ja{もえる} mean? (\textquotedblleft{}To burn\textquotedblright{}.) Say it once more.
\end{tcolorbox}
\section[\texorpdfstring{\ja{きえる} (kieru)}{きえる (kieru)}]{\ja{きえる} (kieru) — to vanish; to go out (of a fire or light)}

\label{lesson:JA-C214-kieru}



\begin{quote}
Before the new one: say the Japanese for a fire, then the Japanese for to burn.
\end{quote}

\subsection*{You'll want to know: \ja{きえる}}
\textbf{\ja{きえる}} — \emph{kieru} — \textquotedblleft{}to vanish; to go out (of a fire or light)\textquotedblright{}.

\begin{quote}
\textbf{\ja{あめで} \ja{ひが} \ja{きえます。}} — \emph{ame de hi ga kiemasu} — \textquotedblleft{}The fire goes out in the rain.\textquotedblright{}
\end{quote}

Polite \textbf{\ja{きえます}}, \emph{kiemasu}; an ichidan verb.

\begin{cousinweb}
\textbf{\ja{きえる}} (\emph{kieru}) is a native Japanese word (wago), inherited rather than borrowed from Chinese; its kanji was attached later as a kun'yomi, and its older origin is not settled (JMdict, KANJIDIC2).
\end{cousinweb}

\begin{grammarlens}[title={What you've built}]
One more word for everyday life at A2.
\end{grammarlens}

\subsection*{Your turn}
\begin{itemize}
\raggedright
  \item \emph{Say it:} \emph{kieru}
  \item \emph{Say it:} \emph{kieru}, once more
  \item \emph{Say it:} \emph{ame de hi ga kiemasu} — The fire goes out in the rain.
  \item \emph{Recall:} say the Japanese for a fire, then the Japanese for to burn, then say \emph{kieru} again
\end{itemize}

\begin{tcolorbox}[breakable,colback=teal!4,colframe=teal!35!black,title={Before you move on}]
What does \ja{きえる} mean? (\textquotedblleft{}To vanish; to go out (of a fire or light)\textquotedblright{}.) Say it once more.
\end{tcolorbox}
